%% toc: P3 40 40.7
\section{Proofs for the jump and complexity conditions}
\label{sec:annex-p03-w03-p03-070}

Throughout, \(\lambda,\mu\in(0,1]\) are reals, graphs are irreflexive and
symmetric, and weights, density, rigidity, safety and closures are those of
\cref{def:p03-w03-p03-070-weight,def:p03-w03-p03-070-closure}.  For a set
\(Y\) of vertices, ``\(Y\) is rigid over \(X\) in \(G\)'' depends only on the
induced subgraph \(G|_{X\cup Y}\), since every edge counted in a weight
\(w_\lambda(X\cup Y_0,Y\setminus Y_0)\) has both ends in \(X\cup Y\).  The
same holds for density and safety.  This is used without further mention
when a statement is transported along an isomorphism of induced subgraphs or
read inside an induced subgraph.

\begin{annexitem}{annex:p03-w03-p03-070-weights}{thm:p03-w03-p03-070-weights}{Calculus of weights}
\begin{proof}
(1) An edge counted in \(e_K(X,Y_1\cup Y_2)\) has both ends in
\(X\cup Y_1\cup Y_2\) and at least one end in \(Y_1\cup Y_2\).  If it has no
end in \(Y_2\), it has both ends in \(X\cup Y_1\) and an end in \(Y_1\), so it
is counted in \(e_K(X,Y_1)\).  Otherwise it has an end in \(Y_2\) and is
counted in \(e_K(X\cup Y_1,Y_2)\).  Conversely, every edge counted in either
term is counted on the left, and no edge is counted in both, since the first
term counts only edges with no end in \(Y_2\).  Hence
\(e_K(X,Y_1\cup Y_2)=e_K(X,Y_1)+e_K(X\cup Y_1,Y_2)\).  Since
\(|Y_1\cup Y_2|=|Y_1|+|Y_2|\), the weights add when the edge sets are finite.
If the edge set on the left is infinite, at least one of the two disjoint
edge sets on the right is infinite.  The corresponding weight and the
weight on the left are both \(-\infty\), so the asserted equality also
holds under the stated convention.

(2) An edge counted in \(e_K(X,Y)\) has an end in \(Y\) and both ends in
\(X\cup Y\subseteq X'\cup Y\), so it is counted in \(e_K(X',Y)\).  Hence
\(e_K(X,Y)\leq e_K(X',Y)\), and since \(\lambda>0\),
\(w_\lambda(X',Y)\leq w_\lambda(X,Y)\).

(3) First suppose that \(X\) is finite.  Since \(Y\) is not safe over
\(X\), the family of nonempty
\(Y_1\subseteq Y\) with \(w_\lambda(X,Y_1)<0\) is nonempty and finite.
Choose \(Y_1\) in it with \(w_\lambda(X,Y_1)\) least and, among those, with
\(|Y_1|\) least.  Let \(Y_0\subsetneq Y_1\).  If \(Y_0=\emptyset\), then
\(Y_1\setminus Y_0=Y_1\) is dense over \(X\).  If \(Y_0\neq\emptyset\), item~(1)
gives
\[
 w_\lambda(X\cup Y_0,Y_1\setminus Y_0)=w_\lambda(X,Y_1)-w_\lambda(X,Y_0).
\]
If \(w_\lambda(X,Y_0)\geq0\), the right side is negative because
\(w_\lambda(X,Y_1)<0\).  If \(w_\lambda(X,Y_0)<0\), then \(Y_0\) lies in the
family, so \(w_\lambda(X,Y_0)\geq w_\lambda(X,Y_1)\).  Equality would
contradict the choice of \(Y_1\) with least size, as \(|Y_0|<|Y_1|\).  Hence
\(w_\lambda(X,Y_0)>w_\lambda(X,Y_1)\), and the right side is again negative.
So \(Y_1\setminus Y_0\) is dense over \(X\cup Y_0\) for every
\(Y_0\subsetneq Y_1\), and \(Y_1\) is rigid over \(X\).

For arbitrary \(X\), choose a nonempty \(D\subseteq Y\) dense over \(X\).
There are more than \(|D|/\lambda\) edges counted in \(e_K(X,D)\), even
when this edge set is infinite.  Choose a finite set of such edges still
numbering more than \(|D|/\lambda\), and let \(X_0\subseteq X\) consist of
their endpoints outside \(D\).  All the chosen edges are counted in
\(e_K(X_0,D)\), so \(D\) is dense over the finite set \(X_0\).  The finite
case just proved gives a nonempty \(Z\subseteq D\) rigid over \(X_0\).
For each proper \(Z_0\subsetneq Z\), monotonicity gives
\[
 w_\lambda(X\cup Z_0,Z\setminus Z_0)
 \leq w_\lambda(X_0\cup Z_0,Z\setminus Z_0)<0.
\]
Thus \(Z\) is rigid over \(X\).  This argument uses only finite weights
in its subtraction step.  Finally, if \(Y\) is dense over \(X\), then
\(Y\) itself is a nonempty subset that is not sparse, so \(Y\) is not safe
and the first statement applies.

(4) Let \(Y'\defeq Y\setminus X'\neq\emptyset\) and \(Y'_0\subsetneq Y'\).  Put
\(Y_0\defeq(Y\cap X')\cup Y'_0\) and \(D\defeq Y'\setminus Y'_0\neq\emptyset\).
Then \(Y\setminus Y_0=D\), so \(Y_0\subsetneq Y\).  Also
\(X\cup Y_0\subseteq X'\cup Y'_0\), and \(D\) is disjoint from
\(X'\cup Y'_0\).  By item~(2) and the rigidity of \(Y\) over \(X\),
\[
 w_\lambda(X'\cup Y'_0,D)\leq w_\lambda(X\cup Y_0,D)<0 .
\]
So \(Y'\setminus Y'_0\) is dense over \(X'\cup Y'_0\) for every
\(Y'_0\subsetneq Y'\), and \(Y'\) is rigid over \(X'\).
\end{proof}
\end{annexitem}

\begin{annexitem}{annex:p03-w03-p03-070-closure}{thm:p03-w03-p03-070-closure}{Closure is well defined}
\begin{proof}
(1) The set \(V(G)\) is \(t\)-closed, since no set lies outside it, so
\(Z\defeq\operatorname{cl}_t(X)\) is an intersection of a nonempty family
\((Z_j)\) of \(t\)-closed sets containing \(X\).  Suppose \(Y\) were a rigid
\(t\)-extension of \(Z\).  For each \(j\), \(Z\subseteq Z_j\), so by
\cref{thm:p03-w03-p03-070-weights}(4), \(Y\setminus Z_j\) is empty or rigid
over \(Z_j\).  In the second case it would be a rigid \(t\)-extension of
\(Z_j\), since \(|Y\setminus Z_j|\leq|Y|\leq t\).  This is impossible, so
\(Y\subseteq Z_j\) for every \(j\), and \(Y\subseteq Z\).  This contradicts
\(Y\cap Z=\emptyset\) and \(Y\neq\emptyset\).  Hence \(Z\) is \(t\)-closed.
It contains \(X\) and lies in every \(t\)-closed set containing \(X\).

(2) By induction on \(j\), \(R_j\subseteq\operatorname{cl}_t(X)\).  For
\(j=0\) this is \(X\subseteq\operatorname{cl}_t(X)\).  If
\(R_j\subseteq\operatorname{cl}_t(X)\), then \(Y_j\defeq R_{j+1}\setminus R_j\)
is rigid over \(R_j\).  By \cref{thm:p03-w03-p03-070-weights}(4),
\(Y_j\setminus\operatorname{cl}_t(X)\) is empty or rigid over
\(\operatorname{cl}_t(X)\) with at most \(t\) elements.  The second case is
excluded by item~(1), so \(R_{j+1}\subseteq\operatorname{cl}_t(X)\).  If
\(R_l\) is \(t\)-closed, it is a \(t\)-closed set containing \(X\), so
\(\operatorname{cl}_t(X)\subseteq R_l\), and equality follows.

(3) A rigid \(t\)-extension has at most \(t\leq t'\) elements, so it is a
rigid \(t'\)-extension.  A set with none of the latter has none of the former.

(4) Since \(W\) is \(t\)-closed and contains \(X\), item~(1) gives
\(\operatorname{cl}^G_t(X)\subseteq W\).  For \(Z\subseteq W\) we claim that
\(Z\) is \(t\)-closed in \(G\) if and only if it is \(t\)-closed in
\(G|_W\).  A rigid \(t\)-extension of \(Z\) in \(G|_W\) is one in \(G\),
since rigidity is read in the induced subgraph.  Conversely, let \(Y\) be a
rigid \(t\)-extension of \(Z\) in \(G\).  By
\cref{thm:p03-w03-p03-070-weights}(4), \(Y\setminus W\) is empty or rigid over
\(W\).  The second case would give a rigid \(t\)-extension of \(W\).  So
\(Y\subseteq W\), and \(Y\) is a rigid \(t\)-extension of \(Z\) in \(G|_W\).
Now \(\operatorname{cl}^G_t(X)\) is a subset of \(W\) that is \(t\)-closed in
\(G\), hence in \(G|_W\), so it contains
\(\operatorname{cl}^{G|_W}_t(X)\).  The latter is \(t\)-closed in \(G|_W\),
hence in \(G\), so it contains \(\operatorname{cl}^G_t(X)\).

(5) A rigid set is dense: take \(Y_0=\emptyset\) in the definition.  So a
rigid \(t\)-extension is a dense set of between \(1\) and \(t\) elements.
Conversely, if \(Y\) is dense over \(X\) with \(1\leq|Y|\leq t\), then
\cref{thm:p03-w03-p03-070-weights}(3) gives a nonempty \(Y_1\subseteq Y\)
rigid over \(X\).  It is disjoint from \(X\) and has at most \(t\) elements,
so it is a rigid \(t\)-extension.
\end{proof}
\end{annexitem}

\begin{annexitem}{annex:p03-w03-p03-070-finite-closure}{thm:p03-w03-p03-070-finite-closure}{Finite Closure Lemma, in counting form}
\begin{proof}
Write \(\varepsilon=\varepsilon_\lambda(t)\) and \(B=\lfloor|X|/\varepsilon\rfloor\).
Run the following process.  Put \(R_0=X\).  While \(R_j\) is not
\(t\)-closed, choose a rigid \(t\)-extension \(Y_j\) of \(R_j\) and put
\(R_{j+1}=R_j\cup Y_j\).  Let \(N_j\defeq|R_j\setminus X|\), so
\(N_0=0\) and \(N_{j+1}=N_j+|Y_j|\) with \(1\leq|Y_j|\leq t\).

Each \(Y_j\) is rigid, hence dense, over \(R_j\).  So
\(v_j\defeq|Y_j|\) and \(e_j\defeq e_G(R_j,Y_j)\) satisfy \(e_j\lambda>v_j\)
with \(1\leq v_j\leq t\).  By the definition of \(\varepsilon\),
\(e_j\lambda-v_j\geq\varepsilon v_j\), that is,
\(w_\lambda(R_j,Y_j)\leq-\varepsilon v_j\).  Repeated use of
\cref{thm:p03-w03-p03-070-weights}(1) gives
\[
 w_\lambda(\emptyset,R_j)=w_\lambda(\emptyset,X)+\sum_{k<j}w_\lambda(R_k,Y_k)
 \leq|X|-\varepsilon N_j ,
\]
where the first step also holds for \(X=\emptyset\) with
\(w_\lambda(\emptyset,\emptyset)=0\), and
\(w_\lambda(\emptyset,X)\leq|X|\) because \(\lambda\geq0\).

Suppose the process reached a stage \(j\) with \(N_j>B\), and take the first
such \(j\).  Then \(N_{j-1}\leq B\), so \(N_j\leq B+t\) and
\(|R_j|\leq|X|+B+t\).  Since \(N_j\geq B+1>|X|/\varepsilon\),
\(w_\lambda(\emptyset,R_j)<0\).  This contradicts the hypothesis applied to
\(F=R_j\).  So \(N_j\leq B\) at every stage.  As \(N_j\) strictly increases
while the process runs, the process stops at some \(R_l\) that is
\(t\)-closed.  By \cref{thm:p03-w03-p03-070-closure}(2),
\(R_l=\operatorname{cl}_t(X)\), and
\(|\operatorname{cl}_t(X)\setminus X|=N_l\leq B\).
\end{proof}
\end{annexitem}

\begin{annexitem}{annex:p03-w03-p03-070-same-theory}{thm:p03-w03-p03-070-same-theory}{The two presentations of \(T_\alpha\) agree}
\begin{proof}
(1) By \cref{def:p03-w03-p03-064-rigid-extension}, \(Y\) is rigid over \(X\)
when \(|Y\setminus S|-e_G(S,Y)\,\alpha<0\) for every \(S\) with
\(X\subseteq S\subsetneq X\cup Y\), where \(e_G(S,Y)\) is the number of edges
with one end in \(Y\setminus S\) and the other in \(X\cup Y\).  Since
\(X\cap Y=\emptyset\), the map \(Y_0\mapsto X\cup Y_0\) is a bijection from
the proper subsets \(Y_0\subsetneq Y\) onto these sets \(S\), and
\(Y\setminus S=Y\setminus Y_0\).  An edge has one end in \(Y\setminus Y_0\)
and the other in \(X\cup Y\) exactly when it has at least one end in
\(Y\setminus Y_0\) and both ends in \((X\cup Y_0)\cup(Y\setminus Y_0)=X\cup Y\).
So \(e_G(S,Y)=e_G(X\cup Y_0,Y\setminus Y_0)\), and the inequality for \(S\)
says that \(Y\setminus Y_0\) is \(\alpha\)-dense over \(X\cup Y_0\).  Taken
over all \(Y_0\subsetneq Y\), these are the conditions of \(\alpha\)-rigidity
in \cref{def:p03-w03-p03-070-weight}.  The second equivalence is item~1 of
\cref{thm:p03-w03-p03-064-rigid-basic}, which applies because \(X\) is
finite.

(2) Let \((R,H)\) be a rooted graph and \(Y=V(H)\setminus R\).  The map
\(S\mapsto S\setminus R\) is a bijection from the sets \(S\) with
\(R\subsetneq S\subseteq V(H)\) onto the nonempty \(Y_1\subseteq Y\).  The
subextension \((R,H|_S)\) has \(|S\setminus R|=|Y_1|\) nonroots, and its edges
not joining two roots are the edges of \(H\) with both ends in \(R\cup Y_1\)
and at least one end in \(Y_1\); there are \(e_H(R,Y_1)\) of them.  So
\((R,H|_S)\) is sparse in the sense of
\cref{def:p03-w03-p03-065-dense-sparse} exactly when
\(w_\alpha(R,Y_1)>0\), while \(\alpha\)-safety asks for
\(w_\alpha(R,Y_1)\geq0\).  These agree, because \(w_\alpha(R,Y_1)\neq0\): if
\(e_H(R,Y_1)=0\) it equals \(|Y_1|\geq1\), and otherwise it vanishes only if
\(\alpha=|Y_1|/e_H(R,Y_1)\), which is rational.

(3) By \cref{def:p03-w03-p03-064-generic}, \(\bar y\) is \(t\)-generic over
\(\bar x\) when, for all distinct \(z_1,\ldots,z_w\) with \(w\leq t\),
distinct from the \(x_i\) and \(y_j\), that form a rigid extension over
\(x_1,\ldots,x_r,y_1,\ldots,y_v\), no \(z_k\) is adjacent to any \(y_j\).  A
rigid extension has at least one nonroot (\cref{def:p03-w03-p03-065-rooted-graph}),
so \(w\geq1\).  Let \(X\) be the set of entries of \(\bar x\bar y\) and
\(Z=\{z_1,\ldots,z_w\}\).  By item~(1), \(z_1,\ldots,z_w\) form a rigid
extension over \(\bar x\bar y\) exactly when \(Z\) is \(\alpha\)-rigid over
\(X\).  So the condition says that no set of between \(1\) and \(t\)
vertices, disjoint from \(X\) and \(\alpha\)-rigid over \(X\), contains a
vertex adjacent to an entry of \(\bar y\).  That is \(t\)-genericity in
\cref{def:p03-w03-p03-070-theory}.

(4) Both theories contain the graph axioms.  The Nonexistence schemas range
over the same graphs \(F\) and assert the same thing, the absence of an
injective map \(V(F)\to V(G)\) carrying edges to edges.  Both Generic
Extension schemas range over the rooted graphs of
\cref{def:p03-w03-p03-065-rooted-graph} that are safe, which by item~(2) is
the same set in both readings, and over all \(t\geq0\).  Fix such
\((R,H)\) and \(t\), a graph \(G\) and distinct \(x_1,\ldots,x_r\).  For a
tuple \(y_1,\ldots,y_v\) compare the conditions.
\begin{itemize}
\item Condition~1 of \cref{def:p03-w03-p03-064-theory-alpha} says that
      \(\bar y\) is an \((R,H)\) extension over \(\bar x\): the \(y_j\) are
      distinct and distinct from the \(x_i\), and every edge of \(H\) with an
      end among the nonroots is present
      (\cref{def:p03-w03-p03-065-rooted-graph}).  This is the distinctness
      requirement together with clause~(1) of
      \cref{def:p03-w03-p03-070-theory}.
\item Condition~2 and clause~(2) are the same requirement.
\item For \(t=0\), condition~3 holds for every tuple by
      \cref{def:p03-w03-p03-064-generic}, and clause~(3) is imposed only for
      \(t\geq1\).  For \(t\geq1\) they agree by item~(3).
\end{itemize}
So, for every \(\bar x\), the same tuples \(\bar y\) witness the two axioms,
and \(G\) satisfies one axiom exactly when it satisfies the other.  A graph is
therefore a model of one theory exactly when it is a model of the other, and
by the completeness theorem the two theories have the same consequences.
\end{proof}
\end{annexitem}

\begin{annexitem}{annex:p03-w03-p03-070-as-holds}{thm:p03-w03-p03-070-as-holds}{Every sentence of \(T_\alpha\) holds almost surely}
\begin{proof}
Let \(\phi\) be a sentence of \(T_\alpha\) in the sense of
\cref{def:p03-w03-p03-070-theory}.  If \(\phi\) is a graph axiom, it holds in
every outcome of \(G(n,n^{-\alpha})\), which is a graph.  Otherwise \(\phi\)
is the Nonexistence or Generic Extension axiom for some datum, and by
\cref{thm:p03-w03-p03-070-same-theory}(4) the axiom \(\phi'\) of
\cref{def:p03-w03-p03-064-theory-alpha} for the same datum satisfies
\(G\entails\phi\iff G\entails\phi'\) for every graph \(G\).  Hence
\(\Pr[G(n,n^{-\alpha})\entails\phi]=\Pr[G(n,n^{-\alpha})\entails\phi']\) for
every \(n\).  By \cref{thm:p03-w03-p03-064-theory-as} the right-hand side tends
to \(1\).  So \(\phi\) holds almost surely
(\cref{def:p03-w03-p03-065-almost-sure}).
\end{proof}
\end{annexitem}

\begin{annexitem}{annex:p03-w03-p03-070-generic-closure}{thm:p03-w03-p03-070-generic-closure}{Closures inside a generic extension}
\begin{proof}
Write \(U=O\cup N'\) and \(C'=\operatorname{cl}^{G'|_U}_t(X)\), taken at
\(\mu\).

\emph{\(C'\) is \(t\)-closed in \(G'\).}  Suppose \(Z\) is a rigid
\(t\)-extension of \(C'\) in \(G'\).  If \(Z\subseteq U\), then \(Z\) is
rigid over \(C'\) in \(G'|_U\), contradicting
\cref{thm:p03-w03-p03-070-closure}(1) in \(G'|_U\).  Otherwise
\(Z_{\mathrm{out}}\defeq Z\setminus U\neq\emptyset\).  Since \(C'\subseteq U\),
\cref{thm:p03-w03-p03-070-weights}(4) shows that \(Z_{\mathrm{out}}\) is rigid
over \(U\).  It has at most \(t\) elements and lies outside \(U\), so by
hypothesis no vertex of \(Z_{\mathrm{out}}\) is adjacent to a vertex of \(N'\).
We show that \(Z_{\mathrm{out}}\) is rigid over \(O\).  Let
\(Z_0\subsetneq Z_{\mathrm{out}}\) and \(D=Z_{\mathrm{out}}\setminus Z_0\).  The
edges counted in \(e_{G'}(U\cup Z_0,D)\) but not in \(e_{G'}(O\cup Z_0,D)\) are
the edges with one end in \(D\) and the other in \(N'\), and there are none.
Hence
\[
 w_\mu(O\cup Z_0,D)=w_\mu(U\cup Z_0,D)<0 ,
\]
the inequality by rigidity over \(U\).  So \(Z_{\mathrm{out}}\) is a rigid
\(t\)-extension of \(O\), and by \cref{thm:p03-w03-p03-070-closure}(3) a rigid
\(T\)-extension of \(O\), since \(t\leq T\).  This contradicts the
\(T\)-closedness of \(O\).

\emph{Conclusion.}  By the first part and
\cref{thm:p03-w03-p03-070-closure}(1) in \(G'\),
\(\operatorname{cl}^{G'}_t(X)\subseteq C'\).  For the reverse inclusion, run the
process of \cref{annex:p03-w03-p03-070-finite-closure} in the finite graph
\(G'|_U\) starting from \(X\).  It stops, because \(U\) is finite, at a set
closed in \(G'|_U\).  By \cref{thm:p03-w03-p03-070-closure}(2) in \(G'|_U\)
that set is \(C'\).  Each step is a rigid \(t\)-extension in \(G'|_U\), hence
in \(G'\).  So \cref{thm:p03-w03-p03-070-closure}(2) in \(G'\) gives
\(C'\subseteq\operatorname{cl}^{G'}_t(X)\).
\end{proof}
\end{annexitem}

\begin{annexitem}{annex:p03-w03-p03-070-extension}{thm:p03-w03-p03-070-extension}{One round of Duplicator's strategy}
\begin{proof}
Write \(\mathrm{BIG}=t_i\) and \(\mathrm{SMALL}=t_{i-1}\), and let
\(X\) be the set of entries of \(\bar x\).  Put
\(\mathrm{OLD}=\operatorname{cl}_{\mathrm{BIG}}(\bar x)\) in \(G\) and
\(\mathrm{OLD}'=\operatorname{cl}_{\mathrm{BIG}}(\bar x')\) in \(G'\).

\emph{Finiteness.}  The Nonexistence schema of \(T^-_\lambda\) forbids every
finite \(F\subseteq V(G)\) with \(w_\lambda(\emptyset,F)<0\), because
\(G|_F\) is a copy of a graph with \(|V|-\lambda|E|<0\).  Since
\(\mathrm{BIG}\geq t_1=1\), \cref{thm:p03-w03-p03-070-finite-closure} gives
\(|\mathrm{OLD}|\leq|X|+\lfloor|X|/\varepsilon_\lambda(\mathrm{BIG})\rfloor\).
The same argument applies in \(G'\).

\emph{Agreement of the notions.}  For a set \(Y\) with \(|Y|\leq t_u\), (A1)
gives \(|Y|<e\lambda\iff|Y|<e\mu\) for every \(e\), so \(Y\) is
\(\lambda\)-dense over \(X\) in a graph exactly when it is \(\mu\)-dense.
Rigidity and safety of a set with at most \(t_u\) elements involve only
subsets of it, so they also agree.  So does \(t\)-closedness for
\(t\leq t_u\).  Hence, for \(t\leq t_u\), the \(t\)-closure of a set in a
finite graph is the same at \(\lambda\) and at \(\mu\): both are the least
\(t\)-closed superset, and the closed sets agree.

\emph{Inside case: \(y\in\mathrm{OLD}\).}  By
\cref{thm:p03-w03-p03-070-closure}(3), \(\mathrm{OLD}\) is
\(\mathrm{SMALL}\)-closed in \(G\).  By
\cref{thm:p03-w03-p03-070-closure}(4),
\(\operatorname{cl}^G_{\mathrm{SMALL}}(\bar xy)=
\operatorname{cl}^{G|_{\mathrm{OLD}}}_{\mathrm{SMALL}}(\bar xy)\).  Put
\(y'=\varphi(y)\).  As \(\varphi\) is an isomorphism of
\(G|_{\mathrm{OLD}}\) onto \(G'|_{\mathrm{OLD}'}\) taking \(\bar xy\) to
\(\bar x'y'\), it maps
\(\operatorname{cl}^{G|_{\mathrm{OLD}}}_{\mathrm{SMALL}}(\bar xy)\) at \(\lambda\)
onto \(\operatorname{cl}^{G'|_{\mathrm{OLD}'}}_{\mathrm{SMALL}}(\bar x'y')\) at
\(\lambda\).  By the agreement of notions, this is the same set at \(\mu\).
By \cref{thm:p03-w03-p03-070-closure}(3),(4) in \(G'\) it equals
\(\operatorname{cl}^{G'}_{\mathrm{SMALL}}(\bar x'y')\).  Let \(\varphi_1\) be
the restriction of \(\varphi\).

\emph{Outside case: \(y\notin\mathrm{OLD}\).}  Let
\(C=\operatorname{cl}^G_{\mathrm{SMALL}}(\bar xy)\) and
\(\mathrm{NEW}=C\setminus\mathrm{OLD}\), so \(y\in\mathrm{NEW}\).  We prove
three claims about \(G\) alone.  They use only \(G\models T^-_\lambda\) and
(A2) for \(\lambda\).

\emph{Claim (a): \(|\mathrm{NEW}|\leq\mathrm{BIG}\).}  If \(i=1\), then
\(\mathrm{SMALL}=0\), \(C=X\cup\{y\}\) and \(\mathrm{NEW}=\{y\}\), since
\(X\subseteq\mathrm{OLD}\).  If \(i\geq2\), let \(X_1=X\cup\{y\}\), so
\(|X_1|\leq s+1\leq u-(i-1)\) and \(\mathrm{SMALL}\geq1\).  By
\cref{thm:p03-w03-p03-070-finite-closure} and (A2) with \(j=i-1\),
\[
 |C\setminus X_1|\leq\Bigl\lfloor\frac{u-(i-1)}{\varepsilon_\lambda(t_{i-1})}\Bigr\rfloor
 \leq t_i-1 .
\]
Since \(\mathrm{NEW}\subseteq(C\setminus X_1)\cup\{y\}\), the claim follows.

\emph{Claim (b): \(\mathrm{NEW}\) is \(\lambda\)-safe over \(\mathrm{OLD}\).}
Otherwise \cref{thm:p03-w03-p03-070-weights}(3) gives a nonempty
\(Y_1\subseteq\mathrm{NEW}\) rigid over \(\mathrm{OLD}\).  By Claim~(a),
\(|Y_1|\leq\mathrm{BIG}\), so \(Y_1\) is a rigid \(\mathrm{BIG}\)-extension
of the \(\mathrm{BIG}\)-closed set \(\mathrm{OLD}\), a contradiction.

\emph{Claim (c): \(C=\operatorname{cl}^{G|_{\mathrm{OLD}\cup\mathrm{NEW}}}_{\mathrm{SMALL}}(\bar xy)\).}
Write \(U=\mathrm{OLD}\cup\mathrm{NEW}\supseteq C\).  The set \(C\) is
\(\mathrm{SMALL}\)-closed in \(G\), hence in \(G|_U\), so the closure in
\(G|_U\) is contained in \(C\).  Conversely, run the process of
\cref{annex:p03-w03-p03-070-finite-closure} in \(G\) from \(X_1\) with
parameter \(\mathrm{SMALL}\).  By the proof there it stops at \(C\), and every
stage lies in \(C\subseteq U\).  Each step is therefore a rigid
\(\mathrm{SMALL}\)-extension in \(G|_U\).  So
\cref{thm:p03-w03-p03-070-closure}(2) in \(G|_U\) puts \(C\) inside the closure
in \(G|_U\).  When \(\mathrm{SMALL}=0\) both sides equal \(X_1\).

\emph{Copying \(\mathrm{NEW}\) into \(G'\).}  Let \(\tilde N\) be a set of new
symbols with a bijection \(\iota:\mathrm{NEW}\to\tilde N\).  Let \(H\) be the
graph on \(\mathrm{OLD}'\cup\tilde N\) that agrees with \(G'\) on
\(\mathrm{OLD}'\) and whose edges meeting \(\tilde N\) are the images under
\(\varphi\cup\iota\) of the edges of \(G|_U\) meeting \(\mathrm{NEW}\).  Then
\(\varphi\cup\iota\) is an isomorphism of \(G|_U\) onto \(H\) taking
\(\mathrm{OLD}\) to \(\mathrm{OLD}'\).  By Claims (a) and (b) and the agreement
of notions, \(\tilde N\) is \(\mu\)-safe over \(\mathrm{OLD}'\) in \(H\).  So
\((\mathrm{OLD}',H)\) is a \(\mu\)-safe rooted graph.  Apply the Generic
Extension axiom of \(T^-_\mu\) for \((\mathrm{OLD}',H)\) with genericity
parameter \(\mathrm{SMALL}\), at the roots listed as the elements of
\(\mathrm{OLD}'\).  It yields a set \(N'\subseteq V(G')\setminus\mathrm{OLD}'\)
and a bijection \(\tilde N\to N'\), such that
\(\psi\defeq\varphi\cup(\text{the composite }\mathrm{NEW}\to\tilde N\to N')\)
is an isomorphism of \(G|_U\) onto \(G'|_{\mathrm{OLD}'\cup N'}\).  Clauses
(1) and (2) fix all pairs meeting \(N'\), and \(\varphi\) fixes the pairs
inside \(\mathrm{OLD}'\).  By clause (3), \(N'\) is
\(\mathrm{SMALL}\)-generic over \(\mathrm{OLD}'\), when
\(\mathrm{SMALL}\geq1\).

Put \(y'=\psi(y)\).  By Claim (c) and the isomorphism, \(\psi(C)\) is the
\(\mathrm{SMALL}\)-closure of \(\bar x'y'\) in
\(G'|_{\mathrm{OLD}'\cup N'}\) at \(\lambda\), hence at \(\mu\).
\Cref{thm:p03-w03-p03-070-generic-closure} applies with \(O=\mathrm{OLD}'\),
which is \(\mathrm{BIG}\)-closed in \(G'\), with \(T=\mathrm{BIG}\) and
\(t=\mathrm{SMALL}\).  For \(\mathrm{SMALL}=0\) its hypothesis is vacuous and
its conclusion is trivial.  It gives
\(\psi(C)=\operatorname{cl}^{G'}_{\mathrm{SMALL}}(\bar x'y')\).  Let
\(\varphi_1=\psi|_C\).  It takes \(x_j\) to \(x'_j\), because
\(x_j\in\mathrm{OLD}\), and \(y\) to \(y'\).

The hypotheses are symmetric in \((G,\lambda)\) and \((G',\mu)\), with
\(\varphi^{-1}\) in place of \(\varphi\), so the last sentence of the lemma
follows.
\end{proof}
\end{annexitem}

\begin{annexitem}{annex:p03-w03-p03-070-transfer}{thm:p03-w03-p03-070-transfer}{Transfer between the theories \(T^-_\lambda\) and \(T^-_\mu\)}
\begin{proof}
Quantifier depth is \(0\) for atomic formulas, is unchanged by negation, is
the maximum for conjunction, and increases by one under a quantifier.
Universal quantifiers are read as \(\neg\exists\neg\).  We prove the
equivalence by induction on \(\theta\), for all \(i\), \(s\), tuples and
\(\varphi\) as in the hypothesis at once.

\emph{Atomic.}  If \(\theta\) is \(z_j=z_k\), then \(x_j=x_k\) iff
\(x'_j=x'_k\), because \(\varphi\) is a bijection with \(\varphi(x_j)=x'_j\).
If \(\theta\) is \(z_j\sim z_k\), the equivalence holds because \(\varphi\) is
an isomorphism of induced subgraphs whose domain contains the entries of
\(\bar x\).  This covers \(i=0\), where \(\operatorname{cl}_0\) is the set of
entries.

\emph{Negation and conjunction.}  These follow at once from the induction
hypothesis.

\emph{Existential.}  Let \(\theta=\exists w\,\psi\) have depth at most \(i\).
Then \(\psi\) has depth at most \(i-1\), so \(i\geq1\).  Renaming the bound
variable does not change depth or meaning, so we may assume \(w=z_{s+1}\) is
none of \(z_1,\ldots,z_s\).  Suppose \(G\entails\theta[\bar x]\), and pick
\(y\) with \(G\entails\psi[\bar xy]\).
\Cref{thm:p03-w03-p03-070-extension} gives \(y'\) and an isomorphism
\(\varphi_1\) of the \(t_{i-1}\)-closures of \(\bar xy\) and \(\bar x'y'\)
taking \(\bar xy\) to \(\bar x'y'\).  Since \(s+1\leq u-(i-1)\), the induction
hypothesis at \(i-1\) gives \(G'\entails\psi[\bar x'y']\), so
\(G'\entails\theta[\bar x']\).  The converse uses the symmetric half of
\cref{thm:p03-w03-p03-070-extension}.

\emph{Sentences.}  Take \(i=u\) and \(s=0\).  By
\cref{thm:p03-w03-p03-070-closure}(5), a rigid \(t_u\)-extension of
\(\emptyset\) would be a finite \(F\) with \(w(\emptyset,F)<0\), which the
Nonexistence schema forbids.  So \(\operatorname{cl}_{t_u}(\emptyset)=\emptyset\)
in both graphs, and the empty map is the required isomorphism.  Hence \(G\)
and \(G'\) agree on every sentence of depth at most \(u\).

\emph{Theories.}  Suppose both theories are consistent, and let \(A\) have
depth at most \(u\) with \(T^-_\lambda\vdash A\).  Fix a model \(G\) of
\(T^-_\lambda\).  Every model \(G'\) of \(T^-_\mu\) agrees with \(G\) on
\(A\), and \(G\entails A\), so \(G'\entails A\).  By the completeness theorem,
\(T^-_\mu\vdash A\).  The converse is symmetric.
\end{proof}
\end{annexitem}

\begin{annexitem}{annex:p03-w03-p03-070-complete}{thm:p03-w03-p03-070-complete}{Completeness of \(T^-_\lambda\) and the zero--one law}
\begin{proof}
(1) For each \(u\geq1\) the strict sequence \(s^\lambda\) is admissible for
\((\lambda,\lambda,u)\).  Condition (A1) is trivial, and (A2) holds by the
definition of \(s^\lambda_{j+1}\).  By
\cref{thm:p03-w03-p03-070-transfer}, any two models of \(T^-_\lambda\)
agree on all sentences of depth at most \(u\), for every \(u\), so they are
elementarily equivalent.  If \(T^-_\lambda\) is consistent and
\(T^-_\lambda\nvdash A\), then \(T^-_\lambda\cup\{\neg A\}\) has a model.
Every model then satisfies \(\neg A\), so \(T^-_\lambda\vdash\neg A\).

(2) Let \(\Phi\subseteq T_\alpha\) be finite.  By
\cref{thm:p03-w03-p03-070-as-holds}, \((\mathrm{AS}_\alpha)\) holds, so
\[
 \Pr\Bigl[G(n,n^{-\alpha})\entails\textstyle\bigwedge\Phi\Bigr]
 \geq1-\sum_{\phi\in\Phi}\Pr\bigl[G(n,n^{-\alpha})\not\entails\phi\bigr]\to1 ,
\]
so for large \(n\) some finite graph satisfies \(\Phi\).  By compactness
\(T_\alpha\) is consistent, and by item~(1) it is complete.  If
\(T_\alpha\vdash A\), then \(A\) follows from a finite \(\Phi\subseteq T_\alpha\),
and \(\Pr[G\entails A]\geq\Pr[G\entails\bigwedge\Phi]\to1\).  Otherwise
completeness gives \(T_\alpha\vdash\neg A\), and the same argument gives
\(\Pr[G\entails\neg A]\to1\), so \(\Pr[G\entails A]\to0\).  In both cases
\(f_A(\alpha)\) exists with the stated value.

(3) Let \(\Phi\subseteq T^-_\lambda\) be finite.  Each sentence of \(\Phi\)
is determined by finitely many answers to questions of the form
``\(v<e\lambda\)?'', with \(v\geq1\) and \(e\geq0\).
\begin{itemize}
\item A Nonexistence axiom for \(F\) asks this of
      \((|V(F)|,|E(F)|)\).
\item A Generic Extension axiom for \((R,H)\) and \(t\) asks it of the pairs
      \((|Y_1|,e_H(R,Y_1))\), \(\emptyset\neq Y_1\subseteq V(H)\setminus R\),
      which decide safety.  Clause~(3) asks it of the pairs that decide,
      for each graph on \(r+v+w\) labelled vertices with \(w\leq t\), whether
      the last \(w\) vertices are rigid over the others.  There are
      finitely many such graphs.
\end{itemize}
Let \(P_\Phi\) be the finite set of pairs involved.  For \((v,e)\in P_\Phi\)
with \(v<e\lambda\), we have \(e\geq1\) and \(v/e<\lambda\).  Let \(\delta>0\)
be smaller than \(\lambda-v/e\) for all these pairs.  For
\(\gamma\in(\lambda-\delta,\lambda)\) and \((v,e)\in P_\Phi\): if
\(v<e\lambda\), then \(\gamma>v/e\), so \(v<e\gamma\).  If \(v\geq e\lambda\),
then \(e\gamma\leq e\lambda\leq v\), so \(v<e\gamma\) fails.  Thus all
answers agree at \(\gamma\) and \(\lambda\).  Each sentence of \(\Phi\) is then
literally the corresponding axiom of \(T^-_\gamma\).  For irrational
\(\gamma\) in this interval, \(T^-_\gamma=T_\gamma\) is consistent by
item~(2), so \(\Phi\) has a model.  By compactness \(T^-_\lambda\) is
consistent, and by item~(1) it is complete.
\end{proof}
\end{annexitem}

\begin{annexitem}{annex:p03-w03-p03-070-xpl}{thm:p03-w03-p03-070-xpl}{Structure of the splitting sets}
\begin{proof}
Write \(t_i(\gamma)\) for the source's sequence and \(s_i(\gamma)\) for the
strict sequence, both for the fixed \(u\).  Two facts are used repeatedly.
First, \(\varepsilon_\gamma(t)\) and \(\varepsilon^{\leq}_\gamma(t)\) do not
increase as \(t\) grows, since the minimum is then over a larger set.  Second,
if no pair satisfies \(v=e\gamma\) with \(v\leq t\), then
\(\varepsilon^{\leq}_\gamma(t)=\varepsilon_\gamma(t)\).  For a rational
\(\gamma=a/b\) in lowest terms, \(v=e\gamma\) with \(v\leq t\) has a solution
exactly when \(a\leq t\), because \(v=ea/b\) forces \(b\mid e\) and then
\(v\geq a\).

(1) Let \(\gamma\in\mathrm{XPL}_u\) have numerator \(a\).  Let \(J\) be the
last index for which \(t_J(\gamma)\) is defined.  For \(j<J\) the step from
\(t_j\) to \(t_{j+1}\) was defined, so \(\varepsilon^{\leq}_\gamma(t_j)>0\), no
equality pair has \(v\leq t_j\), and
\(\varepsilon^{\leq}_\gamma(t_j)=\varepsilon_\gamma(t_j)\).  We show
\(s_j\geq t_j\) for \(j\leq J\) by induction.  This is clear for \(j\leq1\).
If \(s_j\geq t_j\) with \(j<J\), then
\(\varepsilon_\gamma(s_j)\leq\varepsilon_\gamma(t_j)\), so
\[
 \Bigl\lfloor\frac{u-j}{\varepsilon_\gamma(s_j)}\Bigr\rfloor+1
 >\frac{u-j}{\varepsilon_\gamma(s_j)}\geq\frac{u-j}{\varepsilon_\gamma(t_j)} ,
\]
and the integer on the left is at least \(\lceil(u-j)/\varepsilon_\gamma(t_j)\rceil\).
Hence \(s_{j+1}\geq t_{j+1}\).  If the sequence is not well defined, then
\(J\leq u-1\) and \(\varepsilon^{\leq}_\gamma(t_J)=0\), so some \(v\leq t_J\)
has \(v=e\gamma\), and \(a\leq t_J\leq s_J\leq s_u\).  If it is well defined,
then \(J=u\) and \(a\leq t_u\leq s_u\).  In both cases
\(\gamma\in\mathrm{XPL}^\ast_u\).

(2) The numerator of \(1/b\) is \(1\).  For the strict sequence,
\(s_u\geq s_1=1\).  For the source's sequence, either it is not well defined,
or \(t_u\geq t_1=1\).  So \(1/b\in S\), and \(b=1\) gives \(1\in S\).  As
\(1/b\to0\), \(S\) has no least element and \(\inf S=0\).

(3) Let \(L\in(0,1]\).  For \(v\geq1\) let \(e_v=\lfloor v/L\rfloor+1\), the
least \(e\) with \(eL>v\).  For \(t\geq1\), let \(\delta_t\in(0,L/2)\) be
smaller than \(L-v/e\) for each pair with \(1\leq v\leq t\),
\(1\leq e\leq2t/L\) and \(v/e<L\).  There are finitely many such pairs.

\emph{Claim.}  For \(\gamma\in(L-\delta_t,L)\), \(1\leq v\leq t\) and \(e\geq0\),
the conditions \(e\gamma\geq v\), \(e\gamma>v\) and \(eL>v\) are equivalent.
If \(e>2t/L\), then \(e\gamma>eL/2>t\geq v\) and \(eL>v\), so all three hold.
If \(e=0\), all three fail.  Otherwise \(1\leq e\leq2t/L\).  If
\(v/e\geq L\), then \(e\gamma<eL\leq v\), and all three fail.  If
\(v/e<L\), then \(v/e\leq L-\delta_t<\gamma\), and all three hold.

By the claim, for such \(\gamma\) the admissible \(e\) for a given \(v\) are
those with \(e\geq e_v\), with no equality.  So
\[
 \varepsilon^{\leq}_\gamma(t)=\varepsilon_\gamma(t)=h_t(\gamma)\defeq
 \min_{1\leq v\leq t}\frac{e_v\gamma-v}{v} ,
\]
a continuous function with \(h_t(L)=\varepsilon_L(t)>0\).  Choose
\(\delta'_t\in(0,\delta_t]\) with \(h_t(\gamma)\geq\varepsilon_L(t)/2\) on
\((L-\delta'_t,L)\).

We show by induction on \(i\leq u\) that there are \(B_i\geq1\) and
\(\eta_i>0\) such that, for every real \(\gamma\in(L-\eta_i,L)\), the source's
sequence is defined up to \(t_i(\gamma)\), and
\(t_i(\gamma),s_i(\gamma)\leq B_i\).  For \(i=1\) take \(B_1=1\) and
\(\eta_1=L/2\).  Given \(B_i,\eta_i\) with \(i\leq u-1\), put
\[
 \eta_{i+1}=\min\Bigl(\eta_i,\min_{1\leq t\leq B_i}\delta'_t\Bigr),\qquad
 \zeta=\min_{1\leq t\leq B_i}\frac{\varepsilon_L(t)}{2}>0 .
\]
For \(\gamma\in(L-\eta_{i+1},L)\) and \(t=t_i(\gamma)\leq B_i\),
\(\varepsilon^{\leq}_\gamma(t)=h_t(\gamma)\geq\zeta>0\).  So
\(t_{i+1}(\gamma)\) is defined and is at most
\(\max(B_i,\lceil(u-i)/\zeta\rceil)\).  Likewise
\(\varepsilon_\gamma(s_i(\gamma))\geq\zeta\), so
\(s_{i+1}(\gamma)\leq\max(B_i,\lfloor(u-i)/\zeta\rfloor+1)\).  Put
\(B_{i+1}=\max(B_i,\lfloor(u-i)/\zeta\rfloor+1)\).

Now every \(\gamma\in S\cap(L-\eta_u,L)\) is a rational at which both
sequences are defined.  Its numerator is at most \(t_u(\gamma)\) or
\(s_u(\gamma)\), hence at most \(B_u\).  A rational \(a/b\in(L/2,L)\) with
\(a\leq B_u\) has \(b<2a/L\leq2B_u/L\), so there are finitely many.  Taking
\(\delta\leq\eta_u\) below the distance from \(L\) to each of them gives
\(S\cap(L-\delta,L)=\emptyset\).

(4) Let \(L\in(0,1)\setminus S\).  If \(S=\mathrm{XPL}_u\), then either \(L\)
is irrational, or \(L\) is rational, the source's sequence is defined at
\(L\), and the numerator of \(L\) exceeds \(t_u(L)\).  In both cases the
sequence is defined at \(L\).  Put \(M=t_u(L)\).  If
\(S=\mathrm{XPL}^\ast_u\), put \(M=s_u(L)\).  In every case no pair has
\(v=eL\) with \(v\leq M\).  Hence \(\varepsilon^{\leq}_L(t)=\varepsilon_L(t)\)
for \(t\leq M\).

The pairs with \(1\leq v\leq M\) and \(eL<v\) have \(e<M/L\), so there are
finitely many.  Choose \(\delta\in(0,1-L)\) smaller than \(v/e-L\) for each of
them with \(e\geq1\).  For \(\gamma\in(L,L+\delta)\), \(1\leq v\leq M\) and
\(e\geq0\), the conditions \(e\gamma\geq v\), \(e\gamma>v\) and \(eL>v\) are
then equivalent.  If \(eL>v\), then \(e\gamma>eL>v\).  If \(eL<v\), then
\(e=0\) or \(\gamma<v/e\), so \(e\gamma<v\).  The case \(eL=v\) does not
occur.  Hence, with \(e_v\) the least \(e\) with \(eL>v\), for \(t\leq M\),
\[
 \varepsilon^{\leq}_\gamma(t)=\varepsilon_\gamma(t)=\min_{v\leq t}\frac{e_v\gamma-v}{v}
 \geq\min_{v\leq t}\frac{e_vL-v}{v}=\varepsilon_L(t) .
\]

We show by induction that \(t_i(\gamma)\) is defined and
\(t_i(\gamma)\leq t_i(L)\) for \(i\leq u\) and \(\gamma\in(L,L+\delta)\), in
the case \(S=\mathrm{XPL}_u\).  For \(i=1\) both equal \(1\).  If
\(t_i(\gamma)\leq t_i(L)\leq M\) and \(i\leq u-1\), then
\[
 \varepsilon^{\leq}_\gamma(t_i(\gamma))\geq\varepsilon^{\leq}_\gamma(t_i(L))
 \geq\varepsilon_L(t_i(L))=\varepsilon^{\leq}_L(t_i(L))>0 .
\]
So \(t_{i+1}(\gamma)\) is defined, and since \(\lceil\cdot\rceil\) is
nondecreasing, \(t_{i+1}(\gamma)\leq t_{i+1}(L)\).  The same computation with
\(\lfloor\cdot\rfloor+1\) gives \(s_i(\gamma)\leq s_i(L)\) in the case
\(S=\mathrm{XPL}^\ast_u\).

Therefore a point of \(S\cap(L,L+\delta)\) is a rational at which the
relevant sequence is defined, with numerator at most \(M\).  A rational
\(a/b>L\) with \(a\leq M\) has \(b<M/L\), so there are finitely many such
points.  Shrinking \(\delta\) below their distances to \(L\) gives
\(S\cap(L,L+\delta)=\emptyset\).

(5) Let \(\emptyset\neq A\subseteq S\) and \(L=\sup A\in(0,1]\).  If
\(L\notin A\), every interval \((L-\delta,L)\) meets \(A\subseteq S\),
contradicting item~(3).  So \(L\in A\) is its greatest element.

(6) For \(q\in S\), the set \(S\cap(0,q)\) contains \(1/b\) for large \(b\),
so it is nonempty, and by item~(5) it has a greatest element \(q^-\).  By the
choice of \(q^-\), \((q^-,q)\) contains no point of \(S\).

Let \(\lambda\in(0,1)\setminus S\), and let
\(Q=S\cap(\lambda,1]\), which contains \(1\).  Put \(q=\inf Q\).  By
item~(4), \(q>\lambda\).  If \(q\notin Q\), then every interval
\((q,q+\delta)\) meets \(S\), and \(q<1\) because \(1\in Q\).  Also
\(q\notin S\), since \(q>\lambda\) and \(q\notin Q\).  This contradicts
item~(4) at \(q\).  So \(q=\min Q\in S\).  No point of \(S\) lies in
\([\lambda,q)\), and \(q^-<q\) lies in \(S\), so \(q^-<\lambda\).  Hence
\(\lambda\in(q^-,q)\).  If also \(\lambda\in(p^-,p)\) with \(p\in S\), then
\(p\in Q\), so \(p\geq q\).  If \(p>q\), then \(q\in S\cap(0,p)\), so
\(q\leq p^-<\lambda\), which is false.  So \(p=q\).  This proves
disjointness and covering.
\end{proof}
\end{annexitem}

\begin{annexitem}{annex:p03-w03-p03-070-nonconstant}{thm:p03-w03-p03-070-nonconstant}{The sequences need not be constant between consecutive points of \(\mathrm{XPL}_u\)}
\begin{proof}
With \(u=2\), the source's sequence is \(t_0=0\), \(t_1=1\) and
\(t_2=\max(1,\lceil1/\varepsilon^{\leq}_\gamma(1)\rceil)\).  Here
\(\varepsilon^{\leq}_\gamma(1)=e\gamma-1\) for the least \(e\) with
\(e\gamma\geq1\).

For \(\gamma\in[2/5,3/7]\) we have \(1/\gamma\in[7/3,5/2]\), so that least
\(e\) is \(3\), and \(3\gamma>1\).  Hence the sequence is defined and
\(t_2(\gamma)=\lceil1/(3\gamma-1)\rceil\).
\begin{itemize}
\item At \(2/5\): \(3\gamma-1=1/5\), \(t_2=5\), and the numerator \(2\) is at
      most \(5\).  So \(2/5\in\mathrm{XPL}_2\).
\item At \(3/7\): \(3\gamma-1=2/7\), \(t_2=\lceil7/2\rceil=4\), and the
      numerator \(3\) is at most \(4\).  So \(3/7\in\mathrm{XPL}_2\).
\item For \(\gamma\in(2/5,3/7)\): \(3\gamma-1\in(1/5,2/7)\), so
      \(1/(3\gamma-1)\in(7/2,5)\) and \(t_2(\gamma)\leq5\).  A rational in this
      interval is in \(\mathrm{XPL}_2\) only if its numerator \(a\) is at most
      \(5\).  Now \(a/b\in(2/5,3/7)\) means \(b\in(7a/3,5a/2)\).  For
      \(a=1,2,3,4\) these intervals are \((7/3,5/2)\), \((14/3,5)\),
      \((7,15/2)\) and \((28/3,10)\), and contain no integer.  For \(a=5\) the
      interval \((35/3,25/2)\) contains only \(b=12\).  But
      \(t_2(5/12)=\lceil1/(5/4-1)\rceil=4<5\), so \(5/12\notin\mathrm{XPL}_2\).
\end{itemize}
So \(\mathrm{XPL}_2\cap(2/5,3/7)=\emptyset\), and \((3/7)^-=2/5\).

For irrational \(\alpha\in(2/5,5/12)\), \(3\alpha-1\in(1/5,1/4)\) and
\(1/(3\alpha-1)\in(4,5)\), so \(t_2(\alpha)=5\).  For irrational
\(\alpha\in(5/12,3/7)\), \(3\alpha-1\in(1/4,2/7)\) and
\(1/(3\alpha-1)\in(7/2,4)\), so \(t_2(\alpha)=4\).

For the strict sequence at \(5/12\), the least \(e\) with \(5e/12>1\) is
\(3\), so \(\varepsilon_{5/12}(1)=1/4\) and
\(s_2=\max(1,\lfloor4\rfloor+1)=5\).  The numerator \(5\) is at most \(5\),
so \(5/12\in\mathrm{XPL}^\ast_2\).  Finally \(5-12\alpha\) is positive for
\(\alpha<5/12\) and negative for \(\alpha>5/12\).
\end{proof}
\end{annexitem}

\begin{annexitem}{annex:p03-w03-p03-070-local}{thm:p03-w03-p03-070-local}{Local admissibility of the strict sequence}
\begin{proof}
Put \(M=s_u\), where \(s=s^\lambda\).  For \(v\geq1\) let
\(e_v=\lfloor v/\lambda\rfloor+1\), the least \(e\) with \(e\lambda>v\).

\emph{(A1) on the left.}  Let \(\delta_1=\min_{1\leq v\leq M}(\lambda-v/e_v)>0\).
Let \(\gamma\in(\lambda-\delta_1,\lambda)\), \(1\leq v\leq M\) and
\(e\geq0\).  If \(v<e\lambda\), then \(e\geq e_v\), so
\(v/e\leq v/e_v<\gamma\) and \(v<e\gamma\).  If \(v\geq e\lambda\), then
\(e\gamma\leq e\lambda\leq v\).  So (A1) holds for \((\gamma,\lambda)\).

\emph{(A1) on the right, under the extra hypothesis.}  If \(\lambda\) is
irrational, or rational with numerator greater than \(M\), no pair has
\(v=e\lambda\) with \(v\leq M\).  The pairs with \(1\leq v\leq M\) and
\(v>e\lambda\) have \(e<M/\lambda\), so there are finitely many.  Let
\(\delta_2>0\) be smaller than \(v/e-\lambda\) for each of them with
\(e\geq1\).  For \(\gamma\in(\lambda,\lambda+\delta_2)\): if \(v<e\lambda\),
then \(v<e\gamma\); if \(v>e\lambda\), then \(e=0\) or \(\gamma<v/e\), so
\(v>e\gamma\).  So (A1) holds for \((\gamma,\lambda)\).

\emph{(A2).}  For \(\kappa=\lambda\), (A2) is the definition of \(s\).  Let
\(\gamma\) be in an interval on which (A1) holds for \((\gamma,\lambda)\).
For \(t\leq M\) and \(v\leq t\), the admissible \(e\) at \(\gamma\) are then
those with \(e\geq e_v\).  So
\(\varepsilon_\gamma(t)=h_t(\gamma)\defeq\min_{v\leq t}(e_v\gamma-v)/v\).
This function is continuous in \(\gamma\) and positive near \(\lambda\), with
\(h_t(\lambda)=\varepsilon_\lambda(t)\).  For \(1\leq j\leq u-1\), let
\(x_j=(u-j)/\varepsilon_\lambda(s_j)\) and \(F_j=\lfloor x_j\rfloor\), so
\(x_j<F_j+1\).  By continuity there is \(\delta_3>0\) such that
\((u-j)/h_{s_j}(\gamma)<F_j+1\) for all \(j\) and all \(\gamma\) with
\(|\gamma-\lambda|<\delta_3\).  For such \(\gamma\),
\(\lfloor(u-j)/\varepsilon_\gamma(s_j)\rfloor+1\leq F_j+1\leq s_{j+1}\).

Item (1) follows with
\[
 \delta=\min(\delta_1,\delta_3),
\]
and item (2) with \(\delta=\min(\delta_1,\delta_2,\delta_3)\), the interval
being intersected with \((0,1]\).
\end{proof}
\end{annexitem}

\begin{annexitem}{annex:p03-w03-p03-070-local-constancy}{thm:p03-w03-p03-070-local-constancy}{Local constancy and left limits of \(f_A\)}
\begin{proof}
(1) Take \(\delta\) from \cref{thm:p03-w03-p03-070-local}(1), and let
\(\gamma\in(\lambda-\delta,\lambda)\) be irrational.  The strict sequence of
\(\lambda\) is admissible for \((\gamma,\lambda,u)\).
\Cref{thm:p03-w03-p03-070-complete}(2),(3) make \(T_\gamma=T^-_\gamma\) and
\(T^-_\lambda\) consistent.  By \cref{thm:p03-w03-p03-070-transfer},
\(T_\gamma\vdash A\iff T^-_\lambda\vdash A\).  By
\cref{thm:p03-w03-p03-070-complete}(2), \(f_A(\gamma)=1\) exactly when
\(T_\gamma\vdash A\).

(2) The same argument applies with \cref{thm:p03-w03-p03-070-local}(2), for
irrational \(\gamma\) on both sides of \(\lambda\).
\end{proof}
\end{annexitem}

\begin{annexitem}{annex:p03-w03-p03-070-jump}{thm:p03-w03-p03-070-jump}{Jump Condition for the refined splitting set}
\begin{proof}
For \(\lambda\in(0,1]\) put \(F(\lambda)=1\) if \(T^-_\lambda\vdash A\) and
\(F(\lambda)=0\) otherwise.  Let \(q\in\mathrm{XPL}^\ast_u\), and let
\(\alpha<\beta\) be irrationals in \((q^-,q)\).  Put \(J=[\alpha,\beta]\).  By
\cref{thm:p03-w03-p03-070-xpl}(6), \(J\cap\mathrm{XPL}^\ast_u=\emptyset\).

\emph{\(F\) is locally constant on \(J\).}  Let \(\lambda\in J\).  Since
\(\lambda\notin\mathrm{XPL}^\ast_u\),
\cref{thm:p03-w03-p03-070-local-constancy}(2) gives an open interval
\(U_\lambda\ni\lambda\) with \(f_A(\gamma)=F(\lambda)\) for all irrational
\(\gamma\in U_\lambda\).  Let \(\lambda'\in U_\lambda\cap J\).
\Cref{thm:p03-w03-p03-070-local-constancy}(1) gives \(\delta'>0\) with
\(f_A(\gamma)=F(\lambda')\) for irrational \(\gamma\in(\lambda'-\delta',\lambda')\).
The open set \(U_\lambda\cap(\lambda'-\delta',\lambda')\) is nonempty, since
\(U_\lambda\) is an open interval containing \(\lambda'\), so it contains an
irrational \(\gamma\).  Then \(F(\lambda')=f_A(\gamma)=F(\lambda)\).

\emph{Constancy.}  The set \(\{\lambda\in J:F(\lambda)=F(\alpha)\}\) and its
complement in \(J\) are both open in \(J\), by local constancy.  As \(J\) is
connected and the first set contains \(\alpha\), it is all of \(J\).  For
irrational \(\lambda\), \(T^-_\lambda=T_\lambda\), and
\cref{thm:p03-w03-p03-070-complete}(2) gives \(f_A(\lambda)=F(\lambda)\).  So
\(f_A(\alpha)=F(\alpha)=F(\beta)=f_A(\beta)\).

\emph{The value.}  By \cref{thm:p03-w03-p03-070-local-constancy}(1) at
\(\lambda=q\), \(f_A(\gamma)=F(q)\) for irrational \(\gamma\) in some
\((q-\delta,q)\).  This interval meets \((q^-,q)\) in irrationals, so the
constant value is \(F(q)\).

\emph{The last sentence.}  Given \(f=f_A\), let \(u\) be the quantifier depth
of \(A\).  By \cref{thm:p03-w03-p03-070-xpl}(6) the intervals
\((q^-,q)\), \(q\in\mathrm{XPL}^\ast_u\), are the intervals defined by the
splitting set, and on each of them \(f\) is constant.
\end{proof}
\end{annexitem}

\begin{annexitem}{annex:p03-w03-p03-070-spectrum}{thm:p03-w03-p03-070-spectrum}{Spectra lie in the refined splitting set}
\begin{proof}
Let \(q\in(0,1)\) be rational with \(q\notin\mathrm{XPL}^\ast_u\).  By
\cref{thm:p03-w03-p03-070-local-constancy}(2), \(f_A\) is constant on the
irrationals of an interval \((q-\delta,q+\delta)\).  So \(f_A(\beta)\)
converges as \(\beta\to q\) through irrationals, and by
\cref{def:p03-w03-p03-069-spectrum}, \(q\) is not a point of discontinuity.
Hence every point of \(\operatorname{Sp}(A)\) lies in \(\mathrm{XPL}^\ast_u\).
\end{proof}
\end{annexitem}

\begin{annexitem}{annex:p03-w03-p03-070-left-limit}{thm:p03-w03-p03-070-left-limit}{The left limit is truth in \(T^-_\lambda\)}
\begin{proof}
Consistency and completeness of \(T^-_\lambda\) are
\cref{thm:p03-w03-p03-070-complete}(3).  Let \(u\) be the quantifier depth of
\(A\).  By \cref{thm:p03-w03-p03-070-local-constancy}(1), \(f_A\) is
constant, with value \(1\) iff \(T^-_\lambda\vdash A\), on the irrationals of
some \((\lambda-\delta,\lambda)\).  So the limit defining \(g_A(\lambda)\)
exists with that value.

For irrational \(\lambda\), part (2) of the same theorem makes \(f_A\)
constant on the irrationals of a neighbourhood of \(\lambda\), including
\(\lambda\).  So \(g_A(\lambda)=f_A(\lambda)\).

Let \(\alpha\) be irrational.  In the proof of
\cref{thm:p03-w03-p03-070-xpl}(6), the least element \(q\) of
\(\mathrm{XPL}^\ast_u\cap(\alpha,1]\) exists and \(\alpha\in(q^-,q)\).  By
\cref{thm:p03-w03-p03-070-jump}, \(f_A(\alpha)=1\) iff \(T^-_q\vdash A\), that
is, iff \(g_A(q)=1\).
\end{proof}
\end{annexitem}

\begin{annexitem}{annex:p03-w03-p03-070-sizes}{thm:p03-w03-p03-070-sizes}{Size bounds at \(\lambda=a/b\)}
\begin{proof}
Let \(1\leq v\leq t\) and \(e\) satisfy \(ec>v\).  Then \(ea-vb\) is a
positive integer, so
\[
 \frac{ec-v}{v}=\frac{ea-vb}{vb}\geq\frac1{vb}\geq\frac1{tb} .
\]
Hence \(\varepsilon_c(t)\geq1/(tb)\).  Then
\[
 s_{j+1}\leq\max\bigl(s_j,(u-j)s_jb+1\bigr)\leq(ub+1)s_j ,
\]
since \((u-j)s_jb+1\leq ubs_j+s_j\) when \(s_j\geq1\).  Induction from
\(s_1=1\) gives \(s_i\leq(ub+1)^{i-1}\).

For the computation, \(s_1=1\).  Given \(s_j\), for each \(v\leq s_j\) the
least admissible \(e\) is \(e_v=\lfloor vb/a\rfloor+1\), and the quotient is
the fraction \((e_va-vb)/(vb)\).  The minimum over \(v\) is found by comparing
fractions through cross multiplication.  Then
\(\lfloor(u-j)/\varepsilon\rfloor\) is an integer division.  All integers
involved are bounded by a polynomial in \(b\) for fixed \(u\), and there are
at most \(s_j\leq(ub+1)^{u-1}\) values of \(v\) at each of \(u-1\) steps.  So
the time is polynomial in \(b\).
\end{proof}
\end{annexitem}

\begin{annexitem}{annex:p03-w03-p03-070-moves}{thm:p03-w03-p03-070-moves}{Permitted moves are exactly the realizable moves}
\begin{proof}
Let \(\mathrm{OLD}=\operatorname{cl}_{s_i}(\bar x)\), and let \(\pi\) be an
isomorphism of \((G|_{\mathrm{OLD}},\bar x)\) onto \((H,\bar h)\).  As in the
finiteness step of \cref{annex:p03-w03-p03-070-extension},
\cref{thm:p03-w03-p03-070-finite-closure} gives
\(|V(H)|=|\mathrm{OLD}|\leq k+\lfloor k/\varepsilon_c(s_i)\rfloor\).

The strict sequence \(s\) is admissible for \((c,c,u)\).  So Claims (a)--(c)
in \cref{annex:p03-w03-p03-070-extension} apply in \(G\), with
\(\lambda=c\), \(\mathrm{BIG}=s_i\) and \(\mathrm{SMALL}=s_{i-1}\).

\emph{Realizable implies permitted.}  Let \(y\in V(G)\) be such that
\((G,\bar xy)\) realizes \((H',\bar h')\).  If \(y\in\mathrm{OLD}\), then by
\cref{thm:p03-w03-p03-070-closure}(3),(4),
\(\operatorname{cl}^G_{s_{i-1}}(\bar xy)=\operatorname{cl}^{G|_{\mathrm{OLD}}}_{s_{i-1}}(\bar xy)\).
Under \(\pi\) this is \(\operatorname{cl}^H_{s_{i-1}}(\bar h\pi(y))\), so the
Inside clause holds with \(\pi(y)\).  If \(y\notin\mathrm{OLD}\), let
\(C=\operatorname{cl}^G_{s_{i-1}}(\bar xy)\) and \(\mathrm{NEW}=C\setminus\mathrm{OLD}\).
Extend \(\pi\) by a bijection of \(\mathrm{NEW}\) onto a set \(N\) of new
vertices.  Let \(K\) be the image of \(G|_{\mathrm{OLD}\cup\mathrm{NEW}}\).
Then \(K\) contains \(H\) as an induced subgraph, \(y_K=\pi(y)\in N\), and
\(1\leq|N|\leq s_i\) by Claim~(a).  By Claim~(b), \(N\) is \(c\)-safe over
\(V(H)\).  By Claim~(c), \((K|_{\operatorname{cl}^K_{s_{i-1}}(\bar hy_K)},\bar hy_K)\)
is the image of \((G|_C,\bar xy)\), which is isomorphic to \((H',\bar h')\).
So the Outside clause holds.

\emph{Permitted implies realizable.}  Inside: given \(y_H\in V(H)\), put
\(y=\pi^{-1}(y_H)\in\mathrm{OLD}\).  As above,
\(\operatorname{cl}^G_{s_{i-1}}(\bar xy)\) is the preimage under \(\pi\) of
\(\operatorname{cl}^H_{s_{i-1}}(\bar hy_H)\), so \((G,\bar xy)\) realizes
\((H',\bar h')\).

Outside: let \(K\), \(N\) and \(y_K\) be given.  Transport \(K\) along
\(\pi^{-1}\) on \(V(H)\) to a rooted graph with roots \(\mathrm{OLD}\).  It is
\(c\)-safe, since \(N\) is safe over \(V(H)\) in \(K\) and \(N\neq\emptyset\).
The Generic Extension axiom of \(T^-_c\), with genericity parameter
\(s_{i-1}\) and roots the elements of \(\mathrm{OLD}\), gives
\(N'\subseteq V(G)\setminus\mathrm{OLD}\).  It also gives an isomorphism
\(\psi\) of \(K\) onto \(G|_{\mathrm{OLD}\cup N'}\) extending \(\pi^{-1}\),
with \(N'\) \(s_{i-1}\)-generic over \(\mathrm{OLD}\) when \(s_{i-1}\geq1\).
Since \(\mathrm{OLD}\) is \(s_i\)-closed and \(s_{i-1}\leq s_i\),
\cref{thm:p03-w03-p03-070-generic-closure} gives, for \(y=\psi(y_K)\),
\[
 \operatorname{cl}^G_{s_{i-1}}(\bar xy)=\operatorname{cl}^{G|_{\mathrm{OLD}\cup N'}}_{s_{i-1}}(\bar xy)
 =\psi\bigl(\operatorname{cl}^K_{s_{i-1}}(\bar hy_K)\bigr).
\]
So \((G,\bar xy)\) realizes \((H',\bar h')\).
\end{proof}
\end{annexitem}

\begin{annexitem}{annex:p03-w03-p03-070-closure-formula}{thm:p03-w03-p03-070-closure-formula}{Closure types are first-order at every parameter}
\begin{proof}
Closures, density and rigidity are taken at \(c\).  As noted at the start of
this annex, they are read inside induced subgraphs, so an isomorphism of
induced subgraphs carries closures computed in one to closures computed in
the other.

\emph{The side condition.}  Let \(W=\operatorname{cl}^G_t(\bar z)\) and
suppose \(\pi\) is an isomorphism of \((G|_W,\bar z)\) onto \((H,\bar h)\).
By \cref{thm:p03-w03-p03-070-closure}(1), \(W\) is \(t\)-closed in \(G\).  By
item~(4) of the same lemma, applied to the set of entries of \(\bar z\) and
to \(W\), \(\operatorname{cl}^{G|_W}_t(\bar z)=W\).  Transporting along
\(\pi\) gives \(\operatorname{cl}^H_t(\bar h)=V(H)\).  So if
\(\operatorname{cl}^H_t(\bar h)\neq V(H)\), no pair \((G,\bar z)\) has the
property, and \(\theta^t_{H,\bar h}\) is taken to be \(\exists w\,\neg(w=w)\),
which no graph satisfies.

\emph{The formula.}  Now let \(\operatorname{cl}^H_t(\bar h)=V(H)\).  List
the vertices of \(H\) that are not entries of \(\bar h\) as
\(d_1,\ldots,d_m\), and take variables \(w_1,\ldots,w_m\).  For
\(a\in V(H)\) let \(\tau(a)\) be \(z_j\) for the least \(j\) with
\(h_j=a\) if \(a\) is an entry, and \(w_l\) if \(a=d_l\).  Put
\begin{itemize}
\item \(\mathrm{Eq}(\bar z)\): the conjunction of \(z_j=z_l\) over the pairs
      with \(h_j=h_l\), and of \(\neg(z_j=z_l)\) over the pairs with
      \(h_j\neq h_l\);
\item \(\mathrm{Diag}(\bar z,\bar w)\): the conjunction of
      \(\neg(\tau(a)=\tau(b))\) over distinct \(a,b\in V(H)\), of
      \(\tau(a)\sim\tau(b)\) over \(ab\in E(H)\), and of
      \(\neg(\tau(a)\sim\tau(b))\) over distinct \(a,b\) with
      \(ab\notin E(H)\);
\item for \(1\leq q\leq t\), \(\mathcal D_q\): the set of sets \(E\) of
      unordered pairs from \(V(H)\cup\{\star_1,\ldots,\star_q\}\), with
      fresh symbols \(\star_1,\ldots,\star_q\), such that every pair in
      \(E\) contains some \(\star_j\) and \(q-c|E|<0\).  It is finite.  For
      \(E\in\mathcal D_q\) let \(\delta_E(\bar z,\bar w,y_1,\ldots,y_q)\) be
      the conjunction of \(\tau'(a)\sim\tau'(b)\) over \(\{a,b\}\in E\), where
      \(\tau'\) extends \(\tau\) by \(\tau'(\star_j)=y_j\);
\item \(\mathrm{Clo}(\bar z,\bar w)\): the conjunction over \(1\leq q\leq t\)
      of
      \[
        \forall y_1\cdots\forall y_q\Bigl(\bigl(\textstyle\bigwedge_{j<j'}
        \neg(y_j=y_{j'})\wedge\bigwedge_{j,\,a\in V(H)}\neg(y_j=\tau(a))\bigr)
        \to\neg\bigvee_{E\in\mathcal D_q}\delta_E\Bigr),
      \]
      which is the empty conjunction, true in every graph, when \(t=0\).
\end{itemize}
Let \(\theta^t_{H,\bar h}\defeq\mathrm{Eq}(\bar z)\wedge\exists w_1\cdots
\exists w_m\,\bigl(\mathrm{Diag}\wedge\mathrm{Clo}\bigr)\).

\emph{What \(\mathrm{Clo}\) says.}  Let \(\iota:V(H)\to V(G)\) be injective,
let \(Z=\iota(V(H))\), and evaluate \(\tau\) through \(\iota\).  Let
\(y_1,\ldots,y_q\) be distinct vertices outside \(Z\), \(1\leq q\leq t\), and
\(Y=\{y_1,\ldots,y_q\}\).  Let \(E^\ast\) be the set of pairs, written with
\(\star_j\) for \(y_j\) and \(a\) for \(\iota(a)\), that are edges of \(G\)
with at least one end in \(Y\) and both ends in \(Z\cup Y\).  Then
\(|E^\ast|=e_G(Z,Y)\).  If \(Y\) is \(c\)-dense over \(Z\), then
\(E^\ast\in\mathcal D_q\) and \(\delta_{E^\ast}\) holds.  Conversely, if
\(\delta_E\) holds for some \(E\in\mathcal D_q\), then
\(E\subseteq E^\ast\), so \(q-c\,e_G(Z,Y)\leq q-c|E|<0\) and \(Y\) is
\(c\)-dense over \(Z\).  Hence \(\mathrm{Clo}\) holds exactly when no set of
between \(1\) and \(t\) vertices outside \(Z\) is \(c\)-dense over \(Z\),
that is, by \cref{thm:p03-w03-p03-070-closure}(5), exactly when \(Z\) is
\(t\)-closed.

\emph{Forward direction.}  Let \(G\entails\theta^t_{H,\bar h}[\bar z]\) with
witnesses \(\bar w\).  By \(\mathrm{Eq}\), the map \(\iota\) with
\(\iota(h_j)=z_j\) and \(\iota(d_l)=w_l\) is well defined.  By
\(\mathrm{Diag}\), it is injective and an isomorphism of \(H\) onto
\(G|_Z\), \(Z=\iota(V(H))\), with \(\iota(\bar h)=\bar z\).  By
\(\mathrm{Clo}\), \(Z\) is \(t\)-closed.  By
\cref{thm:p03-w03-p03-070-closure}(4),
\(\operatorname{cl}^G_t(\bar z)=\operatorname{cl}^{G|_Z}_t(\bar z)
=\iota\bigl(\operatorname{cl}^H_t(\bar h)\bigr)=\iota(V(H))=Z\).  So
\(\iota\) is an isomorphism of \((H,\bar h)\) onto
\((G|_{\operatorname{cl}^G_t(\bar z)},\bar z)\).

\emph{Backward direction.}  Let \(W=\operatorname{cl}^G_t(\bar z)\) and let
\(\pi\) be an isomorphism of \((G|_W,\bar z)\) onto \((H,\bar h)\).  Since
\(\pi(z_j)=h_j\) and \(\pi\) is injective, \(\mathrm{Eq}\) holds.  Put
\(w_l=\pi^{-1}(d_l)\).  With \(\iota=\pi^{-1}\), \(\mathrm{Diag}\) holds, and
\(\iota(V(H))=W\) is \(t\)-closed by \cref{thm:p03-w03-p03-070-closure}(1),
so \(\mathrm{Clo}\) holds.  Hence \(G\entails\theta^t_{H,\bar h}[\bar z]\).
\end{proof}
\end{annexitem}

\begin{annexitem}{annex:p03-w03-p03-070-dichotomy}{thm:p03-w03-p03-070-dichotomy}{The syntactic dichotomy for moves}
\begin{proof}
(1) Let \(G\) be a model of \(T^-_c\) and \(\bar z\in V(G)^k\) with
\(G\entails\theta_P[\bar z]\).  By
\cref{thm:p03-w03-p03-070-closure-formula} with \(t=s_i\), \((G,\bar z)\)
realizes \(P\) at level \(i\).  By \cref{thm:p03-w03-p03-070-moves}, some
\(y\in V(G)\) makes \((G,\bar zy)\) realize \(P'\) if and only if the move
\(P\to P'\) is permitted.  By \cref{thm:p03-w03-p03-070-closure-formula} with
\(t=s_{i-1}\), \((G,\bar zy)\) realizes \(P'\) exactly when
\(G\entails\theta_{P'}[\bar zy]\).  So, if the move is permitted,
\(G\entails\exists y\,\theta_{P'}[\bar z,y]\), and otherwise
\(G\entails\neg\exists y\,\theta_{P'}[\bar z,y]\).  As \(\bar z\) was arbitrary,
\(G\entails\sigma^+_{P,P'}\) in the first case and
\(G\entails\sigma^-_{P,P'}\) in the second.  As \(G\) was an arbitrary model
of \(T^-_c\), the completeness theorem gives
\(T^-_c\vdash\sigma^+_{P,P'}\) or \(T^-_c\vdash\sigma^-_{P,P'}\)
accordingly.

(2) \(T^-_c\) is consistent and complete by
\cref{thm:p03-w03-p03-070-complete}(3).  If \(P\) is realized in a model
\(G_0\), then \(G_0\entails\exists\bar z\,\theta_P(\bar z)\) by
\cref{thm:p03-w03-p03-070-closure-formula}.  By completeness \(T^-_c\) proves
\(\exists\bar z\,\theta_P(\bar z)\) or its negation, and the negation fails in
\(G_0\), so \(T^-_c\vdash\exists\bar z\,\theta_P(\bar z)\).  Conversely, if
\(T^-_c\) proves this sentence, any model, which exists by consistency,
realizes \(P\).  Now let \(P\) be realized.  If \(T^-_c\) proved both
\(\sigma^+_{P,P'}\) and \(\sigma^-_{P,P'}\), it would prove
\(\forall\bar z\,\neg\theta_P(\bar z)\), contradicting its consistency
together with \(T^-_c\vdash\exists\bar z\,\theta_P(\bar z)\).  By item~(1),
\(T^-_c\) proves \(\sigma^+_{P,P'}\) when the move is permitted and
\(\sigma^-_{P,P'}\) when it is not.  So exactly one of the two is provable,
and \(T^-_c\vdash\sigma^+_{P,P'}\) if and only if the move is permitted.
\end{proof}
\end{annexitem}

\begin{annexitem}{annex:p03-w03-p03-070-truth}{thm:p03-w03-p03-070-truth}{Truth is decided by positions}
\begin{proof}
By induction on \(\theta\), for all \(i\), all realizing pairs and all
positions at once.  Every pair \((G,\bar z)\) with \(|\bar z|\leq u-i\)
realizes exactly one level-\(i\) position up to isomorphism, and
\(\mathrm{Val}_i(\theta,\cdot)\) is invariant under isomorphism of positions.

\emph{Atomic.}  The realizing isomorphism takes \(\bar x\) to \(\bar h\), and
the closure contains the entries of \(\bar x\).  So equality and adjacency
between entries agree in \(G\) and in \(H\).

\emph{Negation and conjunction.}  Immediate.

\emph{Existential.}  Let \(\theta=\exists w\,\psi\), with \(w=z_{k+1}\) after
renaming, so \(i\geq1\) and \(\psi\) has depth at most \(i-1\).  For \(y\in V(G)\)
let \(P_y\) be the level-\((i-1)\) position realized by \((G,\bar xy)\); this
is legitimate since \(k+1\leq u-(i-1)\).  Then
\[
 G\entails\theta[\bar x]\iff\exists y\;G\entails\psi[\bar xy]
 \iff\exists y\;\mathrm{Val}_{i-1}(\psi,P_y) ,
\]
by the induction hypothesis.  By \cref{thm:p03-w03-p03-070-moves}, the
positions \(P_y\) are exactly the positions \(P'\) for which \(P\to P'\) is
permitted.  So the last condition is \(\mathrm{Val}_i(\theta,P)\).

\emph{Sentences.}  By the argument in \cref{annex:p03-w03-p03-070-transfer},
\(\operatorname{cl}_{s_u}(\emptyset)=\emptyset\) in every model of
\(T^-_c\).  So \((G,())\) realizes \((\emptyset,())\) at level \(u\).  If
\(T^-_c\) is consistent, then \(T^-_c\vdash A\) iff every model satisfies
\(A\).  Since the value \(\mathrm{Val}_u(A,(\emptyset,()))\) does not depend on
the model, and a model exists, this holds iff \(\mathrm{Val}_u(A,(\emptyset,()))\)
is true.
\end{proof}
\end{annexitem}

\begin{annexitem}{annex:p03-w03-p03-070-complexity}{thm:p03-w03-p03-070-complexity}{Complexity Condition}
\begin{proof}
\emph{Closure properties of the hierarchy.}  A relation on tuples of strings
is placed in \(\Sigma^p_k\) or \(\Pi^p_k\) by encoding the tuple as one string,
under \cref{def:p04-polynomial-hierarchy}.  The following properties are used,
for \(k\geq1\).
\begin{enumerate}
\item[(P1)] \(\mathsf P\subseteq\Sigma^p_k\cap\Pi^p_k\).  Add quantifiers over
      strings the predicate ignores.
\item[(P2)] A relation is in \(\Pi^p_k\) iff its complement is in
      \(\Sigma^p_k\).  Negate the predicate and dualize each quantifier.
\item[(P3)] \(\Sigma^p_k\) and \(\Pi^p_k\) are closed under \(\wedge\) and
      \(\vee\).  Take prefixes \(\exists y_1\forall y_2\cdots\) over \(R\)
      and \(\exists z_1\forall z_2\cdots\) over \(R'\), with disjoint
      variables.  Then the prefix
      \[
       \exists(y_1,z_1)\,\forall(y_2,z_2)\cdots
      \]
      over \(R\wedge R'\) (respectively \(R\vee R'\)) is equivalent.  Pairs of
      strings are encoded with polynomial overhead.
\item[(P4)] \(\Sigma^p_k\) is closed under polynomially bounded \(\exists\):
      merge the new string into the first existential block.  Dually,
      \(\Pi^p_k\) is closed under polynomially bounded \(\forall\).
\item[(P5)] \(\Sigma^p_k\cup\Pi^p_k\subseteq\Sigma^p_{k+1}\cap\Pi^p_{k+1}\).  Add an ignored
      block in front, or at the end.
\end{enumerate}
By (P2) and (P3), each \(\Sigma^p_k\cap\Pi^p_k\) is closed under Boolean operations.

\emph{Encodings and sizes.}  For an input \(0^a1^b\) with \(1\leq a\leq b\), put
\(c=a/b\).  Let \(m(b)=u+u(ub+1)^{u-1}b\).  By
\cref{thm:p03-w03-p03-070-sizes,thm:p03-w03-p03-070-moves}, every position
realized in a model of \(T^-_c\) has at most \(m(b)\) vertices, since
\(k+\lfloor k/\varepsilon_c(s_i)\rfloor\leq u+us_ib\).  A position is encoded
by an adjacency matrix and a tuple of indices.  A subset of a vertex set is
encoded as a bit vector.  All strings below have length polynomial in \(b\),
and \(s_0,\ldots,s_u\) are computed in polynomial time
(\cref{thm:p03-w03-p03-070-sizes}).

\emph{Basic predicates.}  For a finite graph \(K\) and vertex sets in it:
\begin{itemize}
\item \(\mathrm{Dense}(X,Y,K)\), meaning \(|Y|b<a\,e_K(X,Y)\), is in
      \(\mathsf P\);
\item \(\mathrm{Rigid}(X,Y,K)\iff\forall Y_0\subsetneq Y\;
      \mathrm{Dense}(X\cup Y_0,Y\setminus Y_0,K)\) is in \(\Pi^p_1\);
\item \(\mathrm{Safe}(X,Y,K)\iff\forall Y_1\,(\emptyset\neq Y_1\subseteq Y
      \to\neg\mathrm{Dense}(X,Y_1,K))\) is in \(\Pi^p_1\);
\item \(\mathrm{Closed}_t(Z,K)\iff\forall Y\,(Y\subseteq V(K)\setminus Z\wedge
      1\leq|Y|\leq t\to\neg\mathrm{Dense}(Z,Y,K))\) is in \(\Pi^p_1\).  It
      expresses \(t\)-closedness, by \cref{thm:p03-w03-p03-070-closure}(5).
\end{itemize}
Let \(\mathrm{IsCl}_t(X,Z,K)\) hold iff there is a chain
\(X=R_0\subseteq\cdots\subseteq R_l=Z\) with \(l\leq|V(K)|\) such that
\(\mathrm{Rigid}(R_j,R_{j+1}\setminus R_j,K)\) and
\(1\leq|R_{j+1}\setminus R_j|\leq t\) for all \(j<l\), and
\(\mathrm{Closed}_t(Z,K)\).  By \cref{thm:p03-w03-p03-070-closure}(2), this
implies \(Z=\operatorname{cl}^K_t(X)\).  Conversely, the process of
\cref{annex:p03-w03-p03-070-finite-closure}, run in the finite graph \(K\),
produces such a chain ending at \(\operatorname{cl}^K_t(X)\).  So
\(\mathrm{IsCl}_t(X,Z,K)\iff Z=\operatorname{cl}^K_t(X)\).  The universally
quantified conditions for all \(j<l\) combine into one universal block over
\((j,Y_0)\), and \(\mathrm{Closed}_t\) is another.  By (P3) and (P4),
\(\mathrm{IsCl}\) is in \(\Sigma^p_2\).

\emph{(1) Permissibility.}  The move \(P=(H,\bar h)\to P'=(H',\bar h')\) at
level \(i\) is permitted iff one of two conditions holds.  For Inside: there
exist \(y\in V(H)\), \(Z\subseteq V(H)\) and a bijection
\(\rho:Z\to V(H')\) that is an isomorphism of \(H|_Z\) onto \(H'\) taking
\(\bar hy\) to \(\bar h'\), such that
\(\mathrm{IsCl}_{s_{i-1}}(\bar hy,Z,H)\).  For Outside: there exist a graph
\(K\) on \(V(H)\cup N\) with \(1\leq|N|\leq s_i\) inducing \(H\) on \(V(H)\),
\(y\in N\), \(Z\subseteq V(K)\) and an isomorphism \(\rho\) of \(K|_Z\) onto
\(H'\) taking \(\bar hy\) to \(\bar h'\), such that \(\mathrm{Safe}(V(H),N,K)\)
and \(\mathrm{IsCl}_{s_{i-1}}(\bar hy,Z,K)\).  The witnesses have polynomial
size, since \(|N|\leq s_i\).  The checks on \(\rho\), \(K\) and \(N\) are in
\(\mathsf P\), \(\mathrm{Safe}\) is in \(\Pi^p_1\subseteq\Sigma^p_2\), and
\(\mathrm{IsCl}\) is in \(\Sigma^p_2\).  By (P1), (P3), (P4) and (P5), each
disjunct, and hence the disjunction, is in \(\Sigma^p_2\).  This is item (1).

\emph{(2) The recursion.}  Rename the bound variables of \(A\) so that
different quantifier occurrences bind different variables.  The subformulas
form a finite set.  For a subformula \(\theta\) with \(k\) free variables and
a level \(i\) with \(\mathrm{depth}(\theta)\leq i\) and \(k\leq u-i\), let
\(V_{\theta,i}(x,P)\) hold iff \(x=0^a1^b\) with \(1\leq a\leq b\), \(P\) encodes
a level-\(i\) position with \(k\) designated vertices and at most \(m(b)\)
vertices, and \(\mathrm{Val}'_i(\theta,P)\) is true.  Here \(\mathrm{Val}'\)
is defined like \(\mathrm{Val}\), except that the existential step ranges
over permitted \(P'\) with at most \(m(b)\) vertices.

Consider positions realized in a model of \(T^-_c\).  A model exists by
\cref{thm:p03-w03-p03-070-complete}(3).  Every permitted successor of a realized position is
realized, by \cref{thm:p03-w03-p03-070-moves}, and so has at most \(m(b)\)
vertices.  Hence, by induction on \(\theta\),
\(\mathrm{Val}'_i(\theta,P)=\mathrm{Val}_i(\theta,P)\) for realized \(P\).

We show \(V_{\theta,i}\in\Sigma^p_{i+2}\cap\Pi^p_{i+2}\) by induction on \(\theta\).  Atomic
\(\theta\) gives a relation in \(\mathsf P\subseteq\Sigma^p_{i+2}\cap\Pi^p_{i+2}\).  Negation
and conjunction preserve \(\Sigma^p_{i+2}\cap\Pi^p_{i+2}\).  For \(\theta=\exists w\,\psi\)
at level \(i\geq1\),
\[
 V_{\theta,i}(x,P)\iff\exists P'\;\bigl[\mathrm{Perm}(x,i,P,P')\wedge
 V_{\psi,i-1}(x,P')\bigr]\wedge(\text{format checks}),
\]
with \(P'\) ranging over strings of length polynomial in \(b\).  By item~(1),
\(\mathrm{Perm}\in\Sigma^p_2\subseteq\Sigma^p_{i+1}\), since \(i+1\geq2\).  By
induction, \(V_{\psi,i-1}\in\Sigma^p_{i+1}\cap\Pi^p_{i+1}\subseteq\Sigma^p_{i+1}\).  By
(P3) and (P4) the existential is in \(\Sigma^p_{i+1}\), and by (P5) in
\(\Sigma^p_{i+2}\cap\Pi^p_{i+2}\).  The format checks are in \(\mathsf P\).

Finally, \(x\in L_A\) iff \(x=0^a1^b\) with \(1\leq a\leq b\) and
\(V_{A,u}(x,(\emptyset,()))\).  The position \((\emptyset,())\) is realized in
every model, so this value is \(\mathrm{Val}_u(A,(\emptyset,()))\).  By
\cref{thm:p03-w03-p03-070-truth}, together with the consistency of
\(T^-_{a/b}\), it is true iff \(T^-_{a/b}\vdash A\).  Hence
\(L_A\in\Sigma^p_{u+2}\cap\Pi^p_{u+2}\subseteq\mathrm{PH}\).
\end{proof}
\end{annexitem}
